% IV.H Ảnh hưởng của mật độ: mảng bit hay mảng chỉ số (E7) — phần A viết.
\subsection{Mảng bit hay mảng chỉ số: ảnh hưởng của mật độ (E7)}
\label{sec:e7}

\emph{Giả thuyết H2.} Theo Bảng~\ref{tab:do-phuc-tap}, một tidset dạng mảng chỉ số tốn $4\sigma$ byte và phép giao tốn $O(\sigma)$, còn dạng mảng bit tốn $N/8$ byte và $N/64$ phép AND bất kể $\sigma$. Về bộ nhớ, hai dạng cân bằng tại mật độ $\rho = \sigma/N = 1/32$. Về thời gian, điểm cân bằng có thể khác $1/32$: một phép AND trên word 64 bit rẻ hơn nhiều so với một bước trộn có rẽ nhánh, nhưng phép giao thích nghi của v3 lại dừng sớm được.

\emph{Thiết lập.} So v3 (mảng chỉ số) với v4 (mảng bit) trên ba nguồn: (i)~dữ liệu tổng hợp có mật độ điều khiển được ($N = 50.000$, 30 item, mỗi item có mặt độc lập với xác suất $p$ từ 0,003 đến 0,4, ngưỡng $\sigma_{\min} = \lceil 0{,}8 p^2 N \rceil$ để các tập 2 item là phổ biến); (ii)~5 bộ benchmark ở mức minsup cao, trung bình và thấp; (iii)~Groceries và 12 phương án biến đổi ảnh của Ứng dụng~2. Mật độ của một CSDL được tính bằng support trung bình của các item phổ biến chia cho $N$. Vì cần so tương đối trên nhiều điểm, mỗi điểm được đo trong cùng một tiến trình đã khởi động, lấy thời gian nhỏ nhất trong 3 lần.

\emph{Kết quả} (Hình~\ref{fig:f08}). \emph{Về bộ nhớ,} trên dữ liệu tổng hợp tỉ số bộ nhớ tidset (mảng chỉ số chia mảng bit) vượt 1 tại $\rho \approx 0{,}058$, cùng bậc với ngưỡng lý thuyết $1/32 \approx 0{,}031$. \emph{Về thời gian,} mảng bit nhanh hơn ở mọi mật độ của dữ liệu tổng hợp, kể cả $\rho = 0{,}003$ (nhanh hơn 2 lần); lợi thế tăng dần lên 39 lần ở $\rho = 0{,}4$. Trên dữ liệu thật, mảng bit nhanh hơn 3--47 lần trên ba bộ dày, 11--15 lần trên Groceries ($\rho \approx 0{,}04$--$0{,}06$) và tới 17 lần trên các phương án ảnh. Chỉ ở dữ liệu thưa thật có minsup thấp thì mảng chỉ số mới thắng: Retail 0,1\% ($\rho = 0{,}0036$, mảng bit chậm hơn 1,2 lần), Retail 0,01\% (chậm hơn 4,4 lần) và T10I4D100K 0,01\% (chậm hơn 1,8 lần).

\emph{Giải thích.} Điểm cân bằng bộ nhớ cao hơn $1/32$ một chút vì công thức $4\sigma$ so với $N/8$ chỉ đúng cho tidset của lớp gốc. Tidset của các lớp sâu hơn thưa hơn nhiều (một cặp item có mật độ khoảng $\rho^2$), nên dạng mảng chỉ số gọn hơn so với dự đoán. Về thời gian, mật độ không đủ để quyết định. Trên dữ liệu tổng hợp, ngưỡng được đặt để hầu hết các cặp đều phổ biến, nên phép giao dạng mảng không dừng sớm được và mỗi bước trộn có rẽ nhánh vẫn đắt hơn một phép AND trên 64 bit. Trên Retail ở minsup thấp, hơn 99\% phép giao cho kết quả dưới ngưỡng (Bảng~\ref{tab:b04}); phép giao thích nghi của v3 dừng sau vài phần tử, trong khi mảng bit luôn phải AND đủ $N/64$ word.

\emph{Kết luận.} H2 đúng về bộ nhớ và chỉ đúng một phần về thời gian: điểm cân bằng thời gian phụ thuộc vào tỉ lệ phép giao thất bại, không chỉ vào mật độ. Vì vậy bản lai v6, vốn chuyển biểu diễn theo ngưỡng bộ nhớ $1/32$, tiết kiệm bộ nhớ trên dữ liệu thưa nhưng có thể chậm hơn bản nhanh hơn trong hai biểu diễn tới 4,8 lần (ở dữ liệu tổng hợp $\rho \approx 1/32$). Đây là cơ sở cho đề xuất chọn ngưỡng theo thời gian ở mục~\ref{sec:manh-yeu}.

\begin{figure}[t]
\centering
\includegraphics[width=\columnwidth]{F08_mat_do_bitset_vector}
\caption{Tỉ số thời gian mảng chỉ số (v3) chia mảng bit (v4) theo mật độ tidset $\rho$ (cả hai trục theo thang log). Tỉ số lớn hơn 1: mảng bit nhanh hơn. Đường liền: dữ liệu tổng hợp; các điểm: benchmark (mức minsup cao, trung bình, thấp), Groceries và 12 phương án ảnh của Ứng dụng~2; đường đứt dọc: ngưỡng cân bằng bộ nhớ $1/32$.}
\label{fig:f08}
\end{figure}

\emph{Bộ nhớ theo tiến trình duyệt (E10, giả thuyết H7).} Bài báo quan sát đỉnh bộ nhớ lúc dựng các tập 2 item ở đầu mỗi lớp (Fig.~22). Nhóm ghi lượng tidset đang sống sau mỗi tập được xuất, trên Chess 65\% (dày) và T10I4D100K 0,1\% (thưa, cùng họ dữ liệu IBM với bài báo); hình đầy đủ ở Phụ lục (Hình~\ref{fig:f10}). Với cài đặt của nhóm, H7 \emph{không} đúng: đường bộ nhớ gần như phẳng, đỉnh chỉ bằng 1,42--1,47 lần trung bình trên Chess và 1,01--1,04 lần trên T10I4D100K. Trên T10I4D100K, gần như toàn bộ 3,9\,MiB là tidset của lớp gốc, giữ suốt quá trình duyệt; các lớp con nhỏ không đáng kể. Trên Chess, đỉnh nằm ở tập dài 13 item, tức cuối đường đi sâu nhất, vì tidset của dữ liệu dày co lại chậm (support luôn ít nhất 65\% số giao dịch), nên tổng theo độ sâu lớn nhất ở đường dài nhất. Khác biệt với bài báo đến từ cách duyệt: bài báo dựng xong mọi lớp con của một lớp trước khi đệ quy, còn nhóm đệ quy ngay và chỉ giữ một lớp ở mỗi độ sâu (mục~\ref{sec:thuat-toan}), nên không có lúc nào phải giữ cùng lúc tidset của mọi tập 2 item trong một lớp lớn.
